\subsection{An exact finite formula for the mean region count}
\label{sec:region-transfer}

The gauge identity gives more than an upper bound on the number of strict regions. Their mean can be computed from a triangular matrix whose size depends only on the width. The same matrix treats independent centred nondegenerate Gaussian biases by restoring one extra index.

Fix $n\ge1$, put $q=2^n$, and define the binomial tails
\[
\mathcal B_d(t)=\sum_{j=t}^d\binom dj,\qquad 0\le t\le d.
\]
Define a matrix $M=M_n$, with rows and columns indexed by $0,\ldots,n$, by
\begin{equation}\label{eq:region-transfer-matrix}
M_{r,t}=
\begin{cases}
2^{r-t}\binom n{r-t}\mathcal B_{n-r+t}(t),&0\le t\le r\le n,\\
0,&r<t.
\end{cases}
\end{equation}
Let $U$ be its principal submatrix on $1,\ldots,n$, and let
\[
c_t=\binom{n-1}{t-1}\quad(1\le t\le n),\qquad
d_t=\binom nt\quad(0\le t\le n).
\]
All entries are nonnegative integers.

\begin{theorem}[exact mean strict-region counts]\label{thm:region-transfer}
Consider the square independent Gaussian network of \cref{eq:forward}, with input domain $\R^n$, positive variance parameter $\alpha$, and ReLU at every layer, including the output. For every $L\ge1$, its zero-bias strict-region count satisfies
\begin{equation}\label{eq:region-transfer-zero}
\E R_{n,L}=2q^{1-L}\one^T U^{L-1}c.
\end{equation}
If instead all bias coordinates are independent centred Gaussian variables of positive variances, independent of every weight, then the affine strict-region count satisfies
\begin{equation}\label{eq:region-transfer-affine}
\E R^{(b)}_{n,L}=q^{1-L}\one^T M^{L-1}d.
\end{equation}
The bias variances may differ across coordinates and layers. These counts retain different activation patterns even when the final affine output maps agree.
\end{theorem}

The finite-dimensional rank calculation behind both identities is useful on its own.

\begin{lemma}[rank of a formal stack]\label{lem:formal-stack-rank}
Fix intermediate masks $D_1,\ldots,D_{L-1}$ and selected row sets $S_j\subseteq\{1,\ldots,n\}$ from the formal preactivation matrices $H_j^D$. Write $s_j=|S_j|$ and identify each mask with its set of retained coordinates. Almost surely the rank of all selected stacked rows is $r_1$, where
\begin{equation}\label{eq:formal-rank-backward}
r_L=s_L,\qquad r_j=s_j+\min\{r_{j+1},|D_j\setminus S_j|\}.
\end{equation}
At any fixed depth this holds simultaneously for every formal mask and every row selection.
\end{lemma}

\begin{proof}
Work backwards in coefficient space. Set $V_L=\operatorname{span}\{e_i^T:i\in S_L\}$ and
\[
V_j=\operatorname{span}\{e_i^T:i\in S_j\}+V_{j+1}W_{j+1}D_j.
\]
Then $V_{j+1}$ depends only on $W_{j+2},\ldots,W_L$ and the fixed masks and selections, so it is independent of $W_{j+1}$. Conditional on that future, choose an orthonormal row basis of $V_{j+1}$. Multiplication by $W_{j+1}$ gives a Gaussian matrix with independent nondegenerate columns. Its restriction to the $|D_j\setminus S_j|$ coordinates outside $S_j$ has rank $\min\{\dim V_{j+1},|D_j\setminus S_j|\}$ almost surely. Adding the $s_j$ selected coordinate vectors proves \cref{eq:formal-rank-backward}. The span of the original selected rows is $V_1W_1$; $W_1$ is invertible almost surely and preserves its rank. A finite union gives the simultaneous assertion. No activation or realizability event was conditioned on.
\end{proof}

\begin{proof}[Proof of \cref{thm:region-transfer}]
For a fixed intermediate-mask sequence, let $r_D^0$ be the number of open orthants met by its central preactivation stack. Whitney's characteristic-polynomial formula and Zaslavsky's region formula give
\begin{equation}\label{eq:whitney-central}
r_D^0=\sum_{S_1,\ldots,S_L}(-1)^{s_1+\cdots+s_L-r_1}.
\end{equation}
See Stanley~\cite{stanley}, Theorems~2.4 and~2.5(11). Indexed repetitions of a nonzero hyperplane cause no problem: the alternating sum over all nonempty subsets of its copies is $-1$. If a row is zero, the strict orthant count is zero, and the sum in \cref{eq:whitney-central} also vanishes by pairing selections that differ only in that row. Thus the identity includes every formal mask, without treating the whole ambient space as a proper hyperplane.

Summing \cref{eq:orbit} with functional $1$ cancels the $q$ choices for the final mask, which does not enter the stack:
\begin{equation}\label{eq:region-gauge-sum}
\E R_{n,L}=q^{1-L}\sum_{D_1,\ldots,D_{L-1}}\E r_D^0.
\end{equation}
For a future rank $t$ and a current row selection of size $s$, there are $\binom ns$ choices of selected coordinates, $2^s$ choices for the mask inside them, and $\binom{n-s}j$ choices for $j$ retained coordinates outside them. The backward rank is $s+\min(t,j)$. Consequently \cref{eq:whitney-central,eq:region-gauge-sum} give
\[
\E R_{n,L}=q^{1-L}w^T A^{L-1}v,
\quad v_t=(-1)^t\binom nt,\quad w_r=(-1)^r,
\]
where the signed transfer matrix on $0,\ldots,n$ is
\[
A_{r,t}=\sum_{s=0}^n\sum_{j=0}^{n-s}
(-2)^s\binom ns\binom{n-s}j\,
\one_{\{r=s+\min(t,j)\}}.
\]

A change of polynomial basis removes the signs. Associate a vector with its coefficient polynomial in $x$, and write $f_t(x)$ for column $t$ of $A$. Then $f_t(1)=0$, and for $1\le t\le n$,
\[
f_t(x)-f_{t-1}(x)
=(x-1)x^{t-1}\sum_{s=0}^{n-t}(-2)^s\binom ns
\mathcal B_{n-s}(t)x^s.
\]
Indeed, $\min(t,j)$ differs from $\min(t-1,j)$ exactly when $j\ge t$. In the basis $e_t(x)=(x-1)x^{t-1}$ of the polynomials vanishing at $1$, the matrix therefore has entries
\[
(-2)^{r-t}\binom n{r-t}\mathcal B_{n-r+t}(t),\qquad 1\le t\le r\le n.
\]
Conjugation by $\diag((-1)^t)$ gives $U$. The initial polynomial $(1-x)^n$ has $e_t$-coefficient $(-1)^tc_t$, while $e_t(-1)=2(-1)^t$. Evaluation at $-1$ proves \cref{eq:region-transfer-zero}.

For biases, the formal affine preactivations are $H_j^Ds+a_j$, with $a_1=b_1$ and $a_j=b_j+W_jD_{j-1}a_{j-1}$. Conditional on all weights and a fixed mask sequence, the map from $(b_1,\ldots,b_L)$ to $(a_1,\ldots,a_L)$ is block lower triangular with identity diagonal. Thus all formal offsets have a joint density. They need not be independent.

A selected family of affine hyperplanes has nonempty intersection almost surely exactly when its normal rows are linearly independent: independent rows define an onto linear map, while consistent offsets for dependent rows lie in a proper linear subspace. A zero normal row has a nonzero constant almost surely and contributes no hyperplane. The affine versions of Whitney's and Zaslavsky's formulas therefore count independent subsets of the normal rows, since every consistent selected subset has its size equal to its rank.

Suppose the future contains $t$ independent selected rows. For $0\le s\le n-t$, selecting $s$ current rows preserves independence of all $t+s$ rows precisely when at least $t$ mask coordinates outside the selection are retained, by \cref{lem:formal-stack-rank}. There are
\[
\binom ns\,2^s\mathcal B_{n-s}(t)=M_{t+s,t}
\]
choices; selecting more than $n-t$ current rows cannot preserve independence. At the final block the initial counts are $d_t=\binom nt$. The affine gauge identity from \cref{sec:biases}, with the same cancelled final-mask factor $q$, now proves \cref{eq:region-transfer-affine}.
\end{proof}

\subsection{Depth asymptotics at fixed width}
\label{sec:fixed-width-depth}

Keep $n$ fixed and put $\beta_n=1-2^{-n}$. Let $S_{n,L}=\{F^{(\alpha)}_{n,L}\not\equiv0\}$ denote global survival. A strictly negative final preactivation can leave a nonempty strict region whose output is zero, so survival and strict-region counts are different objects.

\begin{corollary}[fixed-width depth asymptotics]\label{cor:fixed-width-depth}
For the zero-bias network, there are finite positive constants $\kappa_n,\sigma_n$, independent of $\alpha>0$, such that
\begin{equation}\label{eq:depth-ratios}
\E R_{n,L}\sim\kappa_n\beta_n^L,\qquad
\PP(S_{n,L})\sim\sigma_n\beta_n^L.
\end{equation}
Both normalized ratios are nondecreasing. The mean constant is explicit: set $m=2^n-1$, $u_1=1$, and
\begin{equation}\label{eq:depth-constant}
u_r=\frac{\sum_{t=1}^{r-1}U_{r,t}u_t}{m-\mathcal B_n(r)}
\quad(2\le r\le n),\qquad
\kappa_n=\frac{2^{n+1}}{2^n-1}\sum_{r=1}^nu_r.
\end{equation}
For $n\ge2$ the mean also has the remainder estimate
\[
\frac{\E R_{n,L}}{\beta_n^L}
=\kappa_n+O_n\!\left(\left[\frac{\mathcal B_n(2)}{2^n-1}\right]^L\right).
\]
Moreover $\kappa_n\ge n\beta_n^{-2}$ and $\beta_n^{-1}\le\sigma_n\le\kappa_n$. In particular,
\begin{equation}\label{eq:depth-rate}
\lim_{L\to\infty}\frac{\log\E R_{n,L}}L
=\lim_{L\to\infty}\frac{\log\PP(S_{n,L})}L
=\log(1-2^{-n}).
\end{equation}
\end{corollary}

\begin{proof}
The diagonal entries of $U$ are the distinct positive tails $\mathcal B_n(1),\ldots,\mathcal B_n(n)$. Its largest eigenvalue is $m=\mathcal B_n(1)$, with positive right eigenvector $u$ given by \cref{eq:depth-constant}. Every other right eigenvector has first coordinate zero, whereas $c_1=1$. Thus the coefficient of $u$ in the spectral decomposition of $c$ is $1$. Apply \cref{eq:region-transfer-zero}. The other eigenvalues are at most $\mathcal B_n(2)<m$, proving the mean limit and its remainder. At $n=1$ the matrix is $U=(1)$ and the same formula gives $\kappa_1=4$ directly.

For monotonicity, let $N_L$ count strict regions whose final mask is nonzero. In the gauge identity the stack is independent of the choice of final mask, so $\E N_L=\beta_n\E R_{n,L}$. Conditional on the first $L$ layers, each such region has postactivation map $x_L(s)=A s$, where $A=D_LH_L^D\ne0$ because at least one retained coordinate is strictly positive. A fresh Gaussian row $w$ satisfies $wA\ne0$ almost surely. Removing the next layer's proper hyperplanes therefore leaves at least one nonempty open strict subregion in each of the finitely many regions counted by $N_L$. Different old activation patterns remain different after extension, so $R_{n,L+1}\ge N_L$ almost surely. Taking expectations proves monotonicity of the normalized means.

For survival, couple all depths by one infinite independent sequence and put $p_L=\PP(S_{n,L})$. On $S_{n,L}$ choose the first input in a fixed enumeration of $\mathbb Q^n$ with nonzero output. It exists by continuity and is measurable with respect to the first $L$ matrices. Conditional on that past, the next fresh layer preserves its output with probability exactly $\beta_n$, so $p_{L+1}\ge\beta_np_L$. Positive homogeneity gives $F(0)=0$, and \cref{lem:reduction} ensures that a nonzero network has a realized strict region. Hence $p_L\le\E R_{n,L}$, and $p_L/\beta_n^L$ is nondecreasing and bounded above by $\kappa_n$. Since $p_1=1$, its limit lies in $[\beta_n^{-1},\kappa_n]$. The lower bound on $\kappa_n$ follows from \cref{eq:depth-strict} below. Positive rescaling of the weights changes neither signs nor zero-output events, so the constants are independent of $\alpha>0$.
\end{proof}

Some useful finite bounds remain particularly simple:
\begin{align}
n\beta_n^{L-2}\ \le\ \E R_{n,L}\ &\le\ C_{n,L}\beta_n^{L-1},&&L\ge2,\label{eq:depth-strict}\\
\beta_n^{L-1}\ \le\ \PP(S_{n,L})\ &\le\ \min\{1,C_{n,L}\beta_n^L\},&&L\ge1.\label{eq:depth-survival}
\end{align}
For the mean lower bound, take a rank-one first mask and nonzero remaining intermediate masks. There are $n(q-1)^{L-2}$ such sequences. All later formal rows are nonzero multiples of a row of $W_1$ almost surely, so their central orthant count is $q$; \cref{eq:region-gauge-sum} gives the result. For the upper bound, a zero intermediate mask kills the strict count, leaving $(q-1)^{L-1}$ sequences, each with count at most $C_{n,L}$. For the survival upper bound, use the invariant functional $\one_{\{J_D\ne0\}}$ in \cref{prop:orbit}; a zero mask at any layer, including the last, kills the Jacobian, leaving at most $(q-1)^L$ contributors. For its lower bound choose $s=W_1^{-1}e_1$ and propagate the first-layer output $e_1$ through the remaining independent layers.

At width one, $\E R_{1,L}=2^{2-L}$ and $\PP(S_{1,L})=2^{1-L}$, so $(\kappa_1,\sigma_1)=(4,2)$. The next mean constants are $\kappa_2=8$ and $\kappa_3=240/7$. No exact higher-width survival constant, convergence rate for the survival ratio, or uniformity in a growing width is asserted.

Earlier work already establishes death at large depth. Rister and Rubin~\cite{rister2021}, Sections~3--4, give global and fixed-input survival bounds with respective bases $1-2^{-n^2}$ and $\beta_n$ in the zero-bias square setting; their Section~5.3 also treats a finite-data survival ratio as width grows, under a variance assumption. Lu et al.~\cite{lu2020}, Theorem~3.3 and Appendix~C, give two-exponential bounds and a one-ray state analysis for one-dimensional input with a final affine readout. After adjusting that depth convention, their bounds and the witness argument above also imply a finite positive survival-ratio limit in the scalar-input setting. A fixed finite set of $M$ inputs needs only $p_L\le M\beta_n^L$ and the same monotonicity. The finite transfer formula supplies the corresponding geometric control over the entire square input space. Neither qualitative eventual death nor ratio arguments in general are claimed as new.

\subsection{The finite limiting mean with independent biases}

\begin{corollary}\label{cor:affine-region-limit}
Under the nondegenerate Gaussian bias assumptions of \cref{thm:region-transfer}, set $z_0=1$ and
\[
z_r=\frac{\sum_{t=0}^{r-1}M_{r,t}z_t}{2^n-\mathcal B_n(r)}
\quad(1\le r\le n),\qquad
\rho_n=\sum_{r=0}^nz_r.
\]
Then $\rho_n$ is a finite positive rational number and
\[
\E R^{(b)}_{n,L}=\rho_n+O_n(\beta_n^L).
\]
On one infinite independent-layer sequence the geometric open strict cells eventually stop changing almost surely. The number of cells in that eventual partition has expectation $\rho_n$.
\end{corollary}

\begin{proof}
The eigenvalues of $M$ are the distinct tails $\mathcal B_n(0),\ldots,\mathcal B_n(n)$. The largest is $q=2^n$, with right eigenvector $z$; since $d_0=z_0=1$, its spectral coefficient is $1$. All remaining eigenvalues are at most $q-1$, so \cref{eq:region-transfer-affine} gives the limit and remainder.

With nondegenerate independent biases, each new preactivation on an old open cell is almost surely either a nonconstant affine function, whose zero set is the intersection of the cell with a proper affine hyperplane, or a nonzero constant. Thus the strict partitions refine: their counts are nondecreasing, and any change of their geometric open cells increases the count. Monotone convergence and the finite mean limit imply a finite integer limiting count almost surely, so only finitely many refinements occur and its expectation is $\rho_n$.
\end{proof}

The first limits are $\rho_1=3$, $\rho_2=21$ and $\rho_3=279$. The finite expectation is stronger than qualitative stabilization alone. After the first layer, a layer with all weights and biases negative maps the whole nonnegative image to zero; it has probability $2^{-n^2-n}$ independently at each depth. Once this happens, later layers cannot recover input dependence. This elementary dropout mechanism also appears in the discussion of Wang~\cite{wang2022}, Section~4.1. The stabilized objects here are geometric activation cells. Their depth-indexed labels and the constant output may keep changing. The exact affine mean requires nondegenerate biases and must not be obtained by substituting zero variance.

\subsection{The width-two formulas}

\begin{proposition}\label{prop:width-two-mean}
For every $L\ge1$ and $\alpha>0$,
\begin{align}
\E R_{2,L}
&=8\left[\left(\frac34\right)^L-\left(\frac14\right)^L\right],\label{eq:width-two-mean}\\
\E R^{(b)}_{2,L}
&=21-24\left(\frac34\right)^L+4\left(\frac14\right)^L,\label{eq:width-two-affine-mean}
\end{align}
where the second identity uses the bias assumptions of \cref{thm:region-transfer}.
\end{proposition}

\begin{proof}
At width two,
\[
U=\begin{pmatrix}3&0\\4&1\end{pmatrix},\quad
M=\begin{pmatrix}4&0&0\\8&3&0\\4&4&1\end{pmatrix},\quad
c=\begin{pmatrix}1\\1\end{pmatrix},\quad
d=\begin{pmatrix}1\\2\\1\end{pmatrix}.
\]
Evaluating the powers in \cref{eq:region-transfer-zero,eq:region-transfer-affine} gives the two identities, including $L=1$ when both means are $4$.
\end{proof}

The first six zero-bias means are $4,4,3.25,2.5,1.890625,1.421875$. Their decay concerns strict regions; ordinary cells remain when the network is identically zero. With nondegenerate biases, the strict partition instead refines and its expected count increases to $21$. These are activation counts throughout, rather than counts of maximal regions of a single output formula.
